\Needspace{6\baselineskip}
\begin{center}\hypertarget{native-265-article191}{}\pdfbookmark[2]{191.}{article-191}191.\end{center}
Quum evolutio theoriae perturbationum ex elementis datis ab instituto
nostro aliena sit, hic tantummodo ostendendum erit, quomodo orbita approximata
ita corrigi possit, ut cum perturbationibus datis iuncta observationibus
satisfaciat quam proxime. Simplicissime hoc negotium absolvitur per methodum
iis quas in artt.~124, 165, 188 exposuimus analogam. Pro temporibus
omnium observationum, quibus ad hunc finem uti propositum est, et quae
prout res fert esse poterunt vel tres vel quatuor vel plures, computabuntur
ex aequationibus perturbationum harum valores numerici, tum pro longitudinibus
in orbita, tum pro radiis vectoribus, tum pro latitudinibus heliocentricis:
ad hunc calculum argumenta desumentur ex elementis ellipticis approximatis,
quibus perturbationum theoria superstructa erat. Dein ex omnibus observationibus
eligentur duae, pro quibus distantiae a terra ex iisdem elementis
approximatis computabuntur: hae hypothesin primam constituent; hypothesis
secunda et tertia formabuntur, distantiis illis paullulum mutatis. In singulis
dein hypothesibus e duobus locis geocentricis determinabuntur positiones heliocentricae
distantiaeque a Sole; ex illis, postquam latitudines a perturbationibus
purgatae fuerint, deducentur longitudo nodi ascendentis, inclinatio orbitae,
longitudinesque in orbita. In hoc calculo methodus art.~110 aliqua modificatione
opus habet, siquidem ad variationem secularem longitudinis nodi et
inclinationis respicere operae pretium videtur. Scilicet designantibus $\beta,\beta'$ latitudines
heliocentricas a perturbationibus periodicis purgatas; $\lambda,\lambda'$ longitudines
heliocentricas; $\AscendingNode,\AscendingNode+\Delta$ longitudines nodi ascendentis; $i,i+\delta$ inclinationes
orbitae; aequationes in hac forma exhibere conveniet:
\[
\begin{gathered}
\mathrm{tang}\,\beta=\mathrm{tang}\,i\sin(\lambda-\AscendingNode)\\
\frac{\mathrm{tang}\,i}{\mathpunct{.}\mathrm{tang}(i+\delta)}\mathrm{tang}\,\beta'
=\mathrm{tang}\,i\sin(\lambda'-\Delta-\AscendingNode).
\end{gathered}
\]
\SourceLeaf{266}Hic valor ipsius $\frac{\mathrm{tang}\,i}{\mathrm{tang}(i+\delta)}$ omni praecisione necessaria obtinetur, substituendo
pro $i$ valorem approximatum: dein $i$ et $\AscendingNode$ per methodos vulgares erui poterunt.

A duabus porro longitudinibus in orbita, nec non a duobus radiis vectoribus
aggregata perturbationum subtrahentur, ut valores pure elliptici prodeant.
Hic vero etiam effectus, quem variationes seculares positionis perihelii et excentricitatis
in longitudinem in orbita radiumque vectorem exserunt, et qui per
formulas differentiales Sect.~I Libri primi determinandus est, statim cum perturbationibus
periodicis iungendus est, siquidem observationes satis ab invicem
distant, ut illius rationem habere operae pretium videatur. Ex his longitudinibus
in orbita radiisque vectoribus correctis, una cum temporibus respondentibus,
elementa reliqua determinabuntur: tandemque ex his elementis positiones
geocentricae pro omnibus reliquis observationibus calculabuntur. Quibus
cum observatis comparatis, eodem modo quem in art.~188 explicavimus systema
id distantiarum elicietur, ex quo elementa omnibus reliquis observationibus
quam optime satisfacientia demanabunt.
